\pg{436} \emph{necessario}. Nem se encontra na \emph{intelligencia creada}; pois esta não forma as verdades \emph{absolutas} e \emph{necessarias}, descobre-as apenas e recebe-as como regra do seu pensamento. Logo a \emph{razão sufficiente} d'aquellas verdades deve encontrar-se numa \emph{intelligencia absoluta} e \emph{necessaria}; pois ellas são entidades \emph{ideaes}, \emph{absolutas} e \emph{necessarias}. A intelligencia \emph{absoluta} e \emph{necessaria} é a de \textsc{Deus}. Logo existe \textsc{Deus} \footnote{E' certo que, além da ordem \emph{physica}, composta de substancias materiaes, existe uma ordem \emph{logica}, composta de verdades \emph{necessarias}, \emph{absolutas}; por quanto estas verdades não podem ser formadas nem mudadas pela nossa intelligencia. Esta faculdade limita-se a descobril-as, e descobre-as como independentes de toda a experiencia e de todo o pensamento creado. Logo ou devemos negar a verdade do principio \emph{de razão sufficiente}, ou deve existir uma intelligencia \emph{necessaria} e \emph{absoluta}, em que estas verdades encontram o seu fundamento. Ora o principio \emph{de razão sufficiente} é verdadeiro. Logo deve existir uma intelligencia \emph{absoluta} e \emph{necessaria}, que é a de \textsc{Deus}.\par Nem se diga que, neste argumento, passamos da ordem \emph{ideal} para a \emph{real}. Por quanto, as verdades \emph{absolutas} e \emph{necessarias} são \emph{objectivas}; pois, como dissemos, ellas são verdades independentemente do nosso pensamento, não podemos formal-as nem mudal-as, são sempre o que são, embora não pensemos nellas e as contestemos. Sendo \emph{objectivas}, devem encontrar na \emph{realidade} um fundamento. Este fundamento não se encontra nas coisas mutaveis, nem na intelligencia creada. Logo deve encontrar-se num Ente \emph{necessario} e \emph{eterno}, em \textsc{Deus}.}.

13. \textbf{Argumentos physicos}. --- Os argumentos \emph{physicos} basêam-se em que, se existe uma ordem maravilhosa no mundo e nas operações dos agentes \emph{irracionaes}, deve existir um ordenador infinitamente sabio e poderoso.

a) Uma ordem maravilhosa e constante existe entre as muitas e diversas substancias, que compõem o universo. Esta ordem exige uma causa infinitamente sabia e poderosa; pois só essa causa, --- por ser \emph{infinitamente sabia}, podia conhecer a natureza de todas as substancias e o posto de cada uma d'ellas no plano da creação, e escolher as leis mais aptas para a constancia e regularidade da mesma ordem; --- e, por ser \emph{infinitamente poderosa}, podia executar esse plano. A causa infinitamente \emph{sabia} e \emph{poderosa} é \textsc{Deus}. Logo existe \textsc{Deus} \footnote{Não insistimos neste ponto; porque na \emph{Cosmologia geral} provámos que a ordem \emph{physica}, com as suas \emph{leis}, é effeito de uma intelligencia \emph{infinita}.\par \emph{Kant} diz: A ordem \emph{physica} é effeito de uma intelligencia, mas não de uma intelligencia \emph{infinita}. Logo não demonstra a existencia de \textsc{Deus}.\par O que dizemos no texto mostra a falsidade da affirmação de \emph{Kant}. --- Além d'isso, a intelligencia, que estabeleceu a ordem no mundo, só podia\pgn{437}ser a do Ente, que \emph{criou} o proprio mundo. Ora o ente, que \emph{criou} o mundo, deve ter um poder \emph{infinito}. O ente, dotado de poder \emph{infinito}, deve ter uma natureza \emph{infinita} e por isso uma intelligencia \emph{infinita}.\par Os materialistas objectam: Esta ordem e regularidade, que existe em todo o universo e em cada substancia, não é effeito de uma causa intelligente e deliberada, mas é obra do \emph{acaso}.\par Resposta. O \emph{acaso}, sendo \emph{privação} de intelligencia e deliberação, é \emph{nada}. Ora é absurdo dizer que o \emph{nada} produziu a \emph{realidade}. Logo a ordem, que existe no mundo, sendo uma \emph{realidade}, não foi produzida pelo \emph{acaso}, mas por um Ente real, infinitamente sabio e poderoso.}.

b) Não só a ordem, que existe no universo, mas tambem a que existe nas operações de cada agente \emph{irracional}, proclama